%----------------------------------------------------------------------------------------
% EXPONENTIËLE FUNCTIES. Alle opgaven zijn eigen werk.
%
% Opbouw: Start, leerstof (theorie, werkwijze of uitleg, voorbeeld, korte oefening), en op een
% nieuwe pagina de Opdrachten (door elkaar, stijgende moeilijkheid), Koppel en Test jezelf.
% Labels van dit hoofdstuk beginnen met h1-.
%----------------------------------------------------------------------------------------

\mychapter{1}{Exponentiële functies}

\lkr{Leerplan B+S'': LPD 12-17 (exponentiële en logaritmische functies), het deel over machten met rationale exponenten, de exponentiële functie en exponentiële groei. Vergelijkingen waarbij een logaritme nodig is, komen in hoofdstuk 2: hier lossen leerlingen alleen vergelijkingen op waarin beide leden een macht van hetzelfde grondtal zijn. De transformaties van grafieken komen uit de cursus Algemene functies.}

Een blad papier is ongeveer $0{,}1$ mm dik. Vouw je het dubbel, dan verdubbelt de dikte, en bij elke volgende vouw opnieuw. Na tien vouwen is het pakje al meer dan tien centimeter dik. Zo'n groei, waarbij elke stap met dezelfde \textit{factor} vermenigvuldigt, heet exponentieel. Je vindt ze bij intrest, bevolkingen, bacteriën en radioactief verval.

\begin{oefeningen}[Start]
\oef{h1-start-1}

\oef{h1-verken}
\end{oefeningen}

%========================================================================================
\section{Machten met rationale exponenten}

Uit de 2de graad ken je machten met gehele exponenten, zoals $2^3 = 8$ en $2^{-3} = \frac18$. Een gebroken exponent definiëren we zo dat de rekenregels blijven gelden: omdat $\bigl(8^{1/3}\bigr)^3 = 8^1$, moet $8^{1/3}$ de derdemachtswortel van $8$ zijn.

\begin{theorie}
Voor $a > 0$, $m \in \mathbb{Z}$ en $n \in \mathbb{N} \setminus \{0\}$ is
\[\formule{a^{\frac{m}{n}} = \sqrt[n]{a^m}} \qquad \text{en} \qquad \formule{a^{-r} = \frac{1}{a^r}}.\]
\begin{stelling}{Rekenregels voor machten}
Voor $a, b > 0$ en rationale $r, s$:
\[\formule{a^r \cdot a^s = a^{r + s}} \qquad \formule{\frac{a^r}{a^s} = a^{r - s}} \qquad \formule{\bigl(a^r\bigr)^s = a^{rs}} \qquad \formule{(ab)^r = a^r b^r} \qquad \formule{\Bigl(\frac{a}{b}\Bigr)^r = \frac{a^r}{b^r}}\]
\end{stelling}
\end{theorie}

\voorbeeld{}{Bereken zonder rekentoestel: $8^{\frac23}$.\label{vb:h1-macht}}{%
\uitleg{De noemer van de exponent is de wortel, de teller de macht. Neem eerst de wortel: dan blijven de getallen klein.}
\begin{stappenbreed}
8^{\frac23} = \bigl(\sqrt[3]{8}\bigr)^2 & Eerst de derdemachtswortel. \\
= 2^2 = 4 &
\end{stappenbreed}}

\begin{oefeningen}
\oef{h1-k-macht}
\end{oefeningen}

\voorbeeld{}{Schrijf als één macht van $x$ (met $x > 0$): $\dfrac{\sqrt{x} \cdot \sqrt[3]{x}}{x^{-\frac16}}$.\label{vb:h1-rekenregels}}{%
\uitleg{Schrijf elke wortel als macht en tel de exponenten op of trek ze af.}
\begin{stappenbreed}
\frac{x^{\frac12} \cdot x^{\frac13}}{x^{-\frac16}} = x^{\frac12 + \frac13 + \frac16} & Wortels als machten, dan de rekenregels. \\
= x^{1} = x & $\frac36 + \frac26 + \frac16 = 1$.
\end{stappenbreed}}

\begin{oefeningen}
\oef{h1-machten}
\end{oefeningen}

%========================================================================================
\section{De exponentiële functie}

\begin{theorie}
Voor een grondtal $a > 0$ met $a \neq 1$ is $\formule{f(x) = a^x}$ een \textbf{\textcolor{accent}{exponentiële functie}}. Ze heeft
\[\formule{\dom f = \R} \qquad \formule{\ber f = \left]0, +\infty\right[}.\]
Haar grafiek gaat door $(0, 1)$ en $(1, a)$, en heeft de $x$-as als \textbf{\textcolor{accent}{horizontale asymptoot}}. Ze is \textbf{stijgend} als $a > 1$ en \textbf{dalend} als $0 < a < 1$.
\end{theorie}

\naast{%
Een asymptoot is een rechte waar de grafiek steeds dichter bij komt, zonder ze te raken: $2^x$ wordt nooit nul, maar voor $x = -10$ is het al kleiner dan $0{,}001$.

Omdat $\bigl(\frac12\bigr)^x = 2^{-x}$, is de grafiek van $\bigl(\frac12\bigr)^x$ het spiegelbeeld van die van $2^x$ om de $y$-as. Hoe groter het grondtal, hoe steiler de grafiek rechts van de $y$-as.}{%
\begin{tikzpicture}[x=.55cm,y=.55cm,font=\scriptsize]
  \rooster{-3}{1}{3}{-1}{1}{5}
  \kromme{-3}{2.32}{2^(\x)}
  \kromme[dashed]{-2.32}{3}{2^(-\x)}
  \kromme[densely dotted]{-3}{1.46}{3^(\x)}
  \vol{(0,1)}
  \node[right,text=accent] at (2.2,4.6) {$2^x$};
  \node[left,text=accent] at (-2.1,4.6) {$\bigl(\frac12\bigr)^x$};
  \node[left,text=accent] at (1.3,4.9) {$3^x$};
\end{tikzpicture}}

\voorbeeld{}{Schets de grafiek van $g(x) = 2^{x + 1} - 3$. Geef het domein, het bereik en de asymptoot.\label{vb:h1-grafiek}}{%
\naast{%
\uitleg{De grafiek van $2^x$ schuift 1 naar links en 3 naar onder (zie Algemene functies). De asymptoot schuift mee naar onder.}
\begin{stappenvrij}{4.2cm}{3.4cm}
(0, 1) \mapsto (-1, -2) & Het vaste punt. \\
\text{asymptoot } y = -3 & De $x$-as, 3 naar onder. \\
\dom g = \R & \\
\ber g = \left]-3, +\infty\right[ & Boven de asymptoot.
\end{stappenvrij}}{%
\begin{tikzpicture}[x=.5cm,y=.5cm,font=\scriptsize]
  \rooster{-4}{1}{2}{-3}{1}{3}
  \draw[grijs,dashed] (-4,-3) -- (2.4,-3);
  \kromme{-4}{1.32}{2^(\x+1)-3}
  \vol{(-1,-2)}
\end{tikzpicture}}}

\begin{oefeningen}
\oef{h1-k-grafiek}
\end{oefeningen}

Een exponentiële functie is overal stijgend of overal dalend, en dus injectief: \textit{verschillende exponenten geven verschillende machten}. Daarmee los je vergelijkingen op waarin beide leden een macht van hetzelfde grondtal zijn.

\voorbeeld{}{Los op: $2^{x + 1} = \dfrac18$.\label{vb:h1-vergelijking}}{%
\uitleg{Schrijf $\frac18$ als macht van 2. Twee machten van 2 zijn gelijk als hun exponenten gelijk zijn.}
\begin{stappenbreed}
2^{x + 1} = 2^{-3} & $\frac18 = \frac{1}{2^3}$. \\
x + 1 = -3 & Exponenten gelijkstellen. \\
x = -4 & Controle: $2^{-3} = \frac18$.
\end{stappenbreed}}

\begin{werkwijze}
Schrijf beide leden als een macht van hetzelfde grondtal (gebruik de rekenregels). Stel dan de exponenten gelijk: $a^p = a^q \iff p = q$. Lukt het niet om hetzelfde grondtal te vinden, dan heb je een logaritme nodig (hoofdstuk 2).
\end{werkwijze}

\begin{oefeningen}
\oef{h1-vergelijking}
\end{oefeningen}

%========================================================================================
\section{Exponentiële groei en verval}

\begin{theorie}
Een grootheid $N$ verandert \textbf{\textcolor{accent}{exponentieel}} als ze per tijdseenheid met dezelfde \textbf{\textcolor{accent}{groeifactor}} $a$ vermenigvuldigd wordt:
\[\formule{N(t) = N_0 \cdot a^t}\]
met $N_0 = N(0)$ de beginwaarde. Bij een groei van $p$\,\% per tijdseenheid is $\formule{a = 1 + \frac{p}{100}}$, bij een daling van $p$\,\% is $\formule{a = 1 - \frac{p}{100}}$. Is $a > 1$, dan spreken we van groei, is $0 < a < 1$, van verval.

De \textbf{\textcolor{accent}{verdubbelingstijd}} is de tijd waarin $N$ verdubbelt, de \textbf{\textcolor{accent}{halveringstijd}} de tijd waarin $N$ halveert.
\end{theorie}

\voorbeeld{}{Je zet $2000$ euro op een rekening met $3$\,\% intrest per jaar. De intrest wordt elk jaar bijgeteld. Hoeveel staat er na 10 jaar op?\label{vb:h1-intrest}}{%
\uitleg{Elk jaar wordt het bedrag vermenigvuldigd met $1{,}03$: de intrest komt bij het kapitaal en levert het jaar erna zelf intrest op.}
\begin{stappenbreed}
K(t) = 2000 \cdot 1{,}03^t & Beginwaarde 2000, groeifactor $1 + \frac{3}{100}$. \\
K(10) = 2000 \cdot 1{,}03^{10} \approx 2687{,}83 & Na 10 jaar, in euro.
\end{stappenbreed}
Met enkelvoudige intrest (elk jaar $60$ euro) was dat maar $2600$ euro geweest.}

\begin{oefeningen}
\oef{h1-k-intrest}
\end{oefeningen}

Ken je de halveringstijd, dan schrijf je het model met de tijd uitgedrukt in halveringstijden.

\voorbeeld{}{Cafeïne heeft in het lichaam een halveringstijd van ongeveer 5 uur. Iemand neemt 200 mg. Hoeveel is er na 12 uur over? Met hoeveel procent daalt de hoeveelheid per uur?\label{vb:h1-verval}}{%
\uitleg{Na $t$ uur zijn er $\frac{t}{5}$ halveringstijden voorbij. De groeifactor per uur is de factor $\frac12$ voor 5 uur, verdeeld over 5 gelijke stappen: $\bigl(\frac12\bigr)^{\frac15}$.}
\begin{stappenbreed}
C(t) = 200 \cdot \Bigl(\frac12\Bigr)^{\frac{t}{5}} & In mg, $t$ in uur. \\
C(12) = 200 \cdot 0{,}5^{2{,}4} \approx 37{,}9 & Na 12 uur, in mg. \\
a = 0{,}5^{\frac15} \approx 0{,}8706 & Groeifactor per uur. \\
1 - 0{,}8706 = 0{,}1294 & Ongeveer $12{,}9$\,\% minder per uur.
\end{stappenbreed}}

\begin{werkwijze}
Zoek de beginwaarde $N_0$ en de groeifactor $a$ per tijdseenheid (uit een percentage, of uit een verdubbelings- of halveringstijd $T$: dan is $N(t) = N_0 \cdot 2^{t/T}$ of $N_0 \cdot \bigl(\frac12\bigr)^{t/T}$). Een groeifactor voor een andere tijdseenheid vind je met de rekenregels: per $k$ eenheden is hij $a^k$.
\end{werkwijze}

\begin{oefeningen}
\oef{h1-k-verval}

\oef{h1-groei}
\end{oefeningen}

%========================================================================================
\section{Het getal e}

Wat gebeurt er als de intrest vaker wordt bijgeteld? Bij $100$\,\% per jaar, $n$ keer per jaar bijgeteld, groeit $1$ euro in een jaar tot $\bigl(1 + \frac1n\bigr)^n$:

\figuren{\renewcommand{\arraystretch}{1.2}%
\begin{tabular}{c|ccccc}
$n$ & $1$ & $10$ & $100$ & $1000$ & $1\,000\,000$ \\ \hline
$\bigl(1 + \frac1n\bigr)^n$ & $2$ & $2{,}5937$ & $2{,}7048$ & $2{,}7169$ & $2{,}7182\,8$
\end{tabular}}

Het resultaat groeit niet onbeperkt, maar nadert een vast getal.

\begin{theorie}
Het \textbf{\textcolor{accent}{getal van Euler}} is het getal dat $\bigl(1 + \frac1n\bigr)^n$ nadert als $n$ onbeperkt groeit:
\[\formule{\e \approx 2{,}718\,28}.\]
De functie $\formule{f(x) = \e^x}$ heet de \textbf{\textcolor{accent}{natuurlijke exponentiële functie}}.
\end{theorie}

Het getal $\e$ is irrationaal, net als $\pi$. Zijn grafiek ligt tussen die van $2^x$ en $3^x$. Waarom juist dit grondtal \textit{natuurlijk} is, zie je bij de afgeleiden: $\e^x$ is de enige exponentiële functie die haar eigen afgeleide is. Bij continue groei met groeisnelheid $r$ (bijvoorbeeld intrest die voortdurend wordt bijgeteld) is $N(t) = N_0 \cdot \e^{rt}$.

\voorbeeld{}{Een bedrag van $1000$ euro groeit continu met $5$\,\% per jaar: $K(t) = 1000\,\e^{0{,}05\,t}$. Hoeveel is dat na 10 jaar? Vergelijk met $5$\,\% per jaar, jaarlijks bijgeteld.\label{vb:h1-e}}{%
\uitleg{Bereken beide met het rekentoestel (de toets $\e^x$).}
\begin{stappenbreed}
K(10) = 1000\,\e^{0{,}5} \approx 1648{,}72 & Continu. \\
1000 \cdot 1{,}05^{10} \approx 1628{,}89 & Jaarlijks: iets minder.
\end{stappenbreed}}

\begin{oefeningen}
\oef{h1-e}

\oef{h1-k-e}
\end{oefeningen}

%========================================================================================
\opdrachtenkop

\begin{oefeningen}
\needspace{14\baselineskip}
\oef{h1-o1}

\needspace{12\baselineskip}
\oef{h1-o10}

\needspace{13\baselineskip}
\oef{h1-o2}

\needspace{12\baselineskip}
\oef{h1-o12}

\needspace{12\baselineskip}
\oef{h1-o3}

\oef{h1-o15}

\needspace{10\baselineskip}
\oef{h1-o13}

\needspace{16\baselineskip}
\oef{h1-o11}

\oef{h1-o14}

\needspace{12\baselineskip}
\oef{h1-o4}

\oef{h1-o18}

\oef{h1-o22}

\oef{h1-o16}

\oef{h1-o6}

\oef{h1-o21}

\oef{h1-o5}

\oef{h1-o19}

\oef{h1-o23}

\needspace{14\baselineskip}
\oef{h1-o17}

\oef{h1-o7}

\oef{h1-o20}

\needspace{10\baselineskip}
\oef{h1-o26}

\oef{h1-o8}

\oef{h1-o25}

\oef{h1-o27}

\oef{h1-o24}

\oef{h1-o9}
\end{oefeningen}

%----------------------------------------------------------------------------------------
% KOPPEL: verbanden met andere leerstof
%----------------------------------------------------------------------------------------
\needspace{20\baselineskip}
\opdrachtsectie{Koppel}

\lkr{Rijen, Kansrekenen en Integralen: kijk na wat al gezien is. Het pakket legt de volgorde niet vast.}

\begin{oefeningen}
\oef{h1-koppel-1}

\oef{h1-koppel-2}

\oef{h1-koppel-3}
\end{oefeningen}

%----------------------------------------------------------------------------------------
\opdrachtsectie{Test jezelf}

\begin{oefeningen}
\oef{h1-t1}

\oef{h1-t2}

\oef{h1-t3}

\oef{h1-t4}

\oef{h1-t5}

\oef{h1-t6}

\oef{h1-t7}
\end{oefeningen}

\herhaalschema{%
  \fv{oef:h1-t1}{vb:h1-macht}{oef:h1-machten}%
  \fv{oef:h1-t2}{vb:h1-vergelijking}{oef:h1-vergelijking}%
  \fv{oef:h1-t3}{vb:h1-grafiek}{oef:h1-k-grafiek}%
  \fv{oef:h1-t4}{vb:h1-intrest}{oef:h1-k-intrest}%
  \fv{oef:h1-t5}{vb:h1-verval}{oef:h1-k-verval}%
  \fv{oef:h1-t6}{vb:h1-e}{oef:h1-k-e}%
  \fv{oef:h1-t7}{vb:h1-verval}{oef:h1-groei}}

\ifzelfstudie
\section*{Maak je eigen samenvatting}
Sluit dit boek en schrijf in je schrift: wat $a^{m/n}$ en $a^{-r}$ betekenen, de rekenregels voor machten, de eigenschappen van de grafiek van $a^x$ voor $a > 1$ en voor $0 < a < 1$, het groeimodel $N_0 \cdot a^t$ met groeifactor, groeipercentage en halveringstijd, en de definitie van $\e$. Vergelijk daarna met de kaders.
\fi
